\documentclass[a4paper]{article}
\usepackage{amsfonts}
\usepackage{amsmath}
\usepackage{amssymb}
\usepackage{graphicx}     

\begin{document}
  
\noindent \large Auteur: CouldBeMathijs \\
\noindent \large Studierichting: Informatica\\
\noindent \large Oefeningenreeks 3E nr. 4\\

\medskip

\normalsize

Huidig volume = startvolume - weggestroomde\\

$\Rightarrow V(t) = 5000 - 5000\left(1 - \dfrac{t}{40}\right)^2$\\

Stroomsnelheid = verandering in volume over de tijd $= \dfrac{dV}{dt}$\\

$\Rightarrow \dfrac{dV}{dt} = -5000\left(1 - \dfrac{t}{40}\right)\left(- \dfrac{1}{40}\right) =  250\left(1- \dfrac{t}{40}\right)$\\

Dus:\\

$V'(5) = \dfrac{875}{4}$ L/min.\\

$V'(10) = \dfrac{750}{4}$ L/min.\\

$V'(20) = 125$ L/min.\\



\begin{thebibliography}{999}
%\bibitem{a} Referentie
\end{thebibliography}
\end{document} 
